\section{Supporting Nondeterminism}
\label{sec:extension}

The \langname{} language of \autoref{sec:context-typing} is
deterministic: every term reduces to (at most) one value under any
particular environment. This section considers a calculus,
\langNDname{}, that relaxes this restriction, enabling us to, e.g.,
type programs like \label{ex:DisjCon} from
\autoref{sec:overview}. Concretely, \langNDname{} extends \langname{}
with nondeterministic \emph{generators} that can produce multiple
values across different executions:
\begin{align*}
  \primop ::= \quad & \ldots ~|~ \Code{gen}_b ~|~ \ldots
\end{align*}
Unsurprisingly, the wrinkle introduced by nondeterminism is related to
dependence: whereas the reachability properties of a \langname{} term
were fundamentally tied with the environment under which it is
evaluated, those of a \langname{} term may (or may not) be
\emph{independent} of the surrounding context.

Dealing with this new possibility raises two technical
challenges. First, nondeterministic functions like generators can be
invoked multiple times, producing independent capabilities at each
call; our extended type system must account for this \emph{resource
  duplication}.  Second, we need to generalize our type denotation to
accomodate terms whose reachability guarantees are not completely
dictated by their evaluation environment.

% (\S\ref{subsec:nondet-denotation}).  We conclude with case studies
% (\S\ref{subsec:case-studies}) demonstrating the expressiveness of
% the resulting system.

% The nondeterministic choice $e_1 \oplus e_2$ is then expressible as
% $\zif{\Code{gen}_\Code{bool}\Code{()}}{e_1}{e_2}$.  These generators serve as the
% ``impure external environments'' introduced informally in
% Section~\ref{sec:context-typing}: the environment $\Envir_x(e) \doteq
% \zlet{x}{\Code{nat\_gen()}}{e}$ is now a first-class term that
% introduces a nondeterministic binding for $x$ whose capability is
% \emph{independent} of the surrounding context at each invocation.

\subsection{Resource Duplication and Persistence} %
\label{subsec:duplication}

In \langname{}, the reachability properties of a term are entangled
with its surrounding context: evaluation does not introduce new
capabilities. For example, every inhabitant of
$\ort{\Unit}{\top} \ctsarr \urt{\Nat}{\vnu > 0}$ must capture some
variable with an angelic binding that provides the witnesses needed to
construct every positive number:
\begin{align*}
  z{:}\urt{\Nat}{\top} \vdash \zzlam{x}{z + 1} : \ort{\Unit}{\top} \ctsarr
  \urt{\Nat}{\vnu + 1}
\end{align*}\noindent
Two successive calls to a \langname{} function with this signature
\emph{always} yield entangled results, even if the calls have
independent arguments:
\begin{align}
  z{:}\urt{\Unit}{\top} \vdash
  \zlet{f}{\zzlam{x}{z}}{\zlet{y_1}{f\;()}{\zlet{y_2}{f\;()}{y_1 -
  y_2}}} : \urt{\Nat}{\top} \judgfail
  \tag{$\Code{e_{ND}}$}\label{ex:nondet}
\end{align}\noindent
This fact is encoded the typing rules for let-expressions,
\textsc{T-Let} and \textsc{T-LetD}, which jointly enforce the
``entangled xor consumed'' policy. The rule \textsc{T-Let} allows the
function $f$ to be used in both the bound expression and the body of a
\Code{let}-expression, but the bound expression will be entangled with
$f$. Thus, if we tried to use \textsc{T-Let} to establish
\ref{ex:nondet}, we would get stuck when trying to type
$\Code{y_1 - y_2}$:
\begin{align*}
  & \underline{\ldots \ctcomma y_1{:}\urt{\Nat}{\top} \ctcomma
    y_2{:}\urt{\Nat}{\top} \vdash y_1 - y_2 : \urt{\Nat}{\top}}
    \quad \judgfail \\
  &\empty \ctcomma \underline{z{:}\urt{\Unit}{\top} \ctcomma f{:}\ort{\Unit}{\top}
    \ctsarr \urt{\Nat}{\top} \ctcomma y_1{:}\urt{\Nat}{\top} \vdash
    \zlet{y_2}{f\;()}{e} : \tau} \tag{\textsc{T-Let}} \\
  & \empty \ctcomma z{:}\urt{\Unit}{\top} \ctcomma f{:}\ort{\Unit}{\top} \ctsarr \urt{\Nat}{\top} \vdash \zlet{y_1}{f\;()}{\zlet{y_2}{f\;()}{e}} : \tau
\end{align*}\noindent
On the other hand, \textsc{T-LetD} produces an independent
result, but prevents $f$ from reappearing in the body of a
\Code{let}-expression, leaving us unable to type the second
\Code{let}-expression in \ref{ex:nondet}:
\begin{align*}
  &\underline{\empty \ctstar y_1{:}\urt{\Nat}{\top} \vdash \zlet{y_2}{f\;()}{e} : \tau \quad \judgfail}  \tag{\textsc{T-LetD}}
  \\
  &\empty \ctstar (z{:}\urt{\Unit}{\top} \ctcomma f{:}\ort{\Unit}{\top} \ctsarr \urt{\Nat}{\top}) \vdash \zlet{y_1}{f\;()}{\zlet{y_2}{f\;()}{e}} : \tau
\end{align*}\noindent

Successive applications of a nondeterministic generator like
$\Code{nat\_gen()}$, however, can produce genuinely independent
values. In order to accurately capture the behavior of such functions,
\langNDname{} introduces a \emph{persistent} type operator:
\begin{align*}
  \Code{nat\_gen} : \ctpers (\ort{\Unit}{\top} \ctsarr \urt{\Nat}{\top})
\end{align*}
Persistence allows a binding in the type context to be duplicated,
e.g.,
\begin{align*}
  x{:}\ctpers \tau \; \ctsubseteq{\{x\}} \; x{:}\ctpers \tau \ctstar x{:}\ctpers \tau
\end{align*}\noindent
When combined with \textsc{T-LetD}, a persistent resource can be
freely duplicated across multiple function applications without
entangling their results.

\paragraph{Duplicable resource and persistent modality} In other
BI-logics that support the persistent modality~\cite{Iris}, the
definition of the persistent modality is based on a \emph{duplicable}
resource with respect to the product operation of the underlying
BI-algebra. A duplicable resource is idempotent with respect to the
product operation, i.e., $r \times r = r$, allowing us to duplicate
the generator types in the type context.  In our setting,
$r \times r = r$ only holds when $r$ is a \emph{singleton} world
containing exactly one environment. Following this idea, we define the
persistent modality to hold on formulas that do not depend on a world
containing different choices:

\begin{definition}[Persistent Modality] We introduce a persistence
  modality $\chpers P$ to our context logic such that:
  \begin{align*}
    r \models \chpers P \text{ iff } \exists \sigma. \{\sigma\} \sqsubseteq r \land \{\sigma\} \models P
  \end{align*}\noindent
  Note that we do not require $r$ to be singleton over its entire
  variable domain, which would violate Kripke monotonicity: if a
  capability $r$ is singleton, then for $r'$ such that
  $r' \sqsubseteq r$, $r'$ is not always singleton. Instead, we impose
  the more relaxed requirement that ensures that $\{\sigma\}$ will at
  least contains all the variables needed by $P$,
  i.e. $\projectTo{r}{\Code{FV}(P)}$.
\end{definition}

\paragraph{Persistent type} The persistent type
$\Gamma \vdash e : \ctpers \tau$ indicates that the typing of $e$ does
not rely on the capabilities provided by either the bindings in the
type context $\Gamma$ or \emph{the reduction of the term $e$}. To
illustrate, consider the following three functions:\footnote{The
  nondeterministic choice operator $e_1 \oplus e_2$ is syntactic sugar
  for $\zif{\Code{gen}_\Code{bool}\Code{()}}{e_1}{e_2}$.}
\begin{align*}
  &z{:}\urt{\Unit}{\top} \vdash \zzlam{x}{z} : \ort{\Unit}{\top}
    \ctsarr \urt{\Nat}{\top} \judgok \\
  &\emptyset \vdash \zzlam{x}{\Code{nat\_gen()}} : \ort{\Unit}{\top}
    \ctsarr \urt{\Nat}{\top}  \judgok \\
  &\emptyset \vdash \zzlam{x}{0} \oplus
    \zzlam{x}{\Code{positive\_gen()}} : \ort{\Unit}{\top} \ctsarr
    \urt{\Nat}{\top}  \judgok \\
\end{align*}\noindent
Of these three terms, only the second can be assigned the persistent
modality. The first is obviously not persistent since it depends on
the binding $z$ in the type context, while the reachability properties
of the third \emph{depend on how it is reduced}. To see how, consider
the following type judgement, which is invalid under \langNDname{}'s
call-by-value semantics:
\begin{align*}
  \emptyset \vdash \zlet{f}{\zzlam{x}{0} \oplus \zzlam{x}{\Code{positive\_gen()}}}{\zif{f\;() = 0}{f\;()}{0}} : \urt{\Nat}{\top} \quad \judgfail
\end{align*}\noindent
This term can only reach the value $0$ since $f\;()$ can either reduce to $0$ or to any positive number. In other words, the two applications of $f$ are indeed dependent, and $f$ is not persistent.

\begin{definition}[Persistent Type Denotation]
  The denotation of $\ctpers \tau$ is:
  \begin{align*}
    \denot{\ctpers \tau}_\Sigma\;e \doteq \forall x.\denot{\mstep{e}{x}} \chimpl \chpers (\denot{\tau}_{\Sigma,x{:}\eraserf{\tau}}\;x) \quad (\text{fresh $x$})
  \end{align*}\noindent
  Note that if we had defined the type denotation as
  $\chpers (\denot{\tau}_{\Sigma}\;e)$, it would (incorrectly) contain
  inhabitants like
  $\zzlam{x}{0} \oplus \zzlam{x}{\Code{positive\_gen()}}$. To ensure
  that $e$ is also single-valued, we introduce a universally
  quantified fresh variable $x$ that ranges over all possible
  reduction results of $e$; the persistent modality then ensures only
  that $x$ is single-valued. \BD{This last point is a bit confusing --
    why does it matter if $x$ is single-valued? Do we mean
    'persistant' here?}
\end{definition}

\paragraph{Typing rule for persistent types}
The elimination rule for persistent types is derivable from
\textsc{T-Sub} and the fact $\Gamma \vdash \ctpers \tau <: \tau$
always holds.  The introduction rule is shown below:
\begin{align*}
\mprooftr{40mm}
      {$\Gamma \vdash v : \tau$ \quad $\denot{\Gamma} \models \chpers \denot{\Gamma}$}
      {T-Pers}
      {$\Gamma \vdash v : \ctpers \tau$}
\end{align*}\noindent
\BD{Can we show the derived elimination rule here? We have the space.}
In short, \textsc{T-Pers} allows us to conclude that a value is a
persistant term of type $\ctpers \tau$ if we can establish that it has
type $\tau$ using a type context comprised only of persistent facts,
i.e., $\denot{\Gamma} \models \chpers \denot{\Gamma}$. \BD{Just to
  double check, requiring a value $v$ in this rule means there are
  persistant terms we cannot be checked against a persistant type,
  right?}

\subsection{Type Denotation for Nondeterminism} %
\label{subsec:nondet-denotation}

Nondeterminism breaks the single-valued assumption of the deterministic
language in two ways that require adjustments to the type denotation.

\paragraph{Function Types} The current rules for let binding are inconsistent with the type denotation in the nondeterministic setting. Consider the following judgement:
\begin{align*}
  \emptyset \vdash \zlet{x}{\Code{bool\_gen()}}{\zlam{y}{\Bool}{y = x}} : y{:}\urt{\Bool}{\top} \ctsarr \urt{\Bool}{\top}
\end{align*}\noindent
The current type denotation assumes that the parameter type $y$ cannot depend on the choice of $x$, which is hidden inside the generator invocation within the term; thus this function should return all boolean values. However, the let-expression rule (e.g., \textsc{T-Let}) exposes the hidden capability within the nondeterministic term:
\begin{align*}
  x{:}\urt{\Bool}{\top} \vdash \zlam{y}{\Bool}{y = x} : y{:}\urt{\Bool}{\top} \ctsarr \urt{\Bool}{\top}
\end{align*}\noindent
where $y$ can be chosen equal to $x$, so the function can only return $\true$.
The issue here is that a nondeterministic term exposes more capabilities into the environment during evaluation, which is not accounted for by the current type denotation.

\paragraph{Sum Types}
A similar issue arises for sum types.  The judgement:
\begin{align*}
  \emptyset \vdash \Code{bool\_gen()} :
  \ort{\Bool}{\vnu = \true} \ctoplus \ort{\Bool}{\vnu = \false}
\end{align*}
should hold, since $\Code{bool\_gen()}$ nondeterministically produces
either $\true$ or $\false$.  However, the original denotation
$\denot{\tau_1 \ctoplus \tau_2}\;e \doteq \denot{\tau_1}\;e \choplus
\denot{\tau_2}\;e$ requires the \emph{context} to be split across the
two branches, not the term itself.  Under the empty context, there is
no capability to split, so the judgement fails despite being
semantically valid.

\paragraph{Generalized Denotation}
Both issues stem from the same root cause: the original denotation
evaluates $e$ once and reasons about the result, but a
nondeterministic $e$ can reduce to different values in different
executions.  The fix is to universally quantify over all possible
reduction results of $e$, lifting $e$'s nondeterminism into the
type context:
\begin{align*}
  \denot{x{:}\tau_x \ctsarr \tau}_\Sigma\;e
    &\doteq \forall f.\; \denot{f \doteq e}_\Sigma \chimpl
             \denot{\tau_x}_\Sigma\;x \chimpl
             \denot{\tau}_{\Sigma,f{:}\eraserf{\tau_x \ctsarr \tau}, x{:}\eraserf{\tau_x}}\;(f\;x)
  \\\denot{x{:}\tau_x \ctwand \tau}_\Sigma\;e
    &\doteq \forall f.\; \denot{f \doteq e}_\Sigma \chimpl
             \denot{\tau_x}_\Sigma\;x \chwand
             \denot{\tau}_{\Sigma,f{:}\eraserf{\tau_x\ctwand \tau}, x{:}\eraserf{\tau_x}}\;(f\;x)
  \\\denot{\tau_1 \ctoplus \tau_2}_\Sigma\;e
    &\doteq \forall x.\; \denot{x \doteq e}_\Sigma \chimpl
             \denot{\tau_1}_{\Sigma,x{:}\eraserf{\tau_1}}\;x \choplus \denot{\tau_2}_{\Sigma,x{:}\eraserf{\tau_2}}\;x
\end{align*}
In each case, a universally quantified fresh variable ($f$ or $x$)
ranges over all possible reduction results of $e$, and the original
denotation is applied to that variable.  In the deterministic case,
$e$ has a unique reduction result and the quantifier is trivially
eliminated, so the generalized denotation coincides with the original.

For the function-type issue, the generalized denotation successfully captures the newly produced capabilities induced by the nondeterministic term. Thus the previous function-type judgement fails under the generalized denotation:
\begin{align*}
  &\denot{y{:}\urt{\Bool}{\top} \ctsarr \urt{\Bool}{\top}}\;
    (\zlet{x}{\Code{bool\_gen()}}{\zlam{y}{\Bool}{y = x}})
  \\=\;& \forall f.\; \Code{Atom}(f = (\zlam{y}{\Bool}{y = \true}) \lor
          f = (\zlam{y}{\Bool}{y = \false})) \chimpl
          \\& \quad \forall y. \Code{Atom}(y = \true \lor y = \false) \chimpl  \angelic\,\Code{Atom}(f\;y = \true \lor f\;y = \false)
  \quad\judgfail
\end{align*}
For the sum-type issue, the generalized denotation lifts
$\Code{bool\_gen()}$'s nondeterminism into the context via the fresh
variable $x$, allowing the sum operator to partition it across the
two branches as usual:
\begin{align*}
  &\denot{\ort{\Bool}{\vnu = \true} \ctoplus
          \ort{\Bool}{\vnu = \false}}\;\Code{bool\_gen()}
  \\=\;& \forall x.\; \Code{Atom}(x = \true \lor x = \false) \chimpl
          \demonic\,\Code{Atom}(\vnu = \true) \choplus
          \demonic\,\Code{Atom}(\vnu = \false)
  \quad\judgok
\end{align*}
